\DocumentMetadata{
  pdfstandard=A-2u,
  lang=cs
}

\tagpdfsetup{activate-all}
\documentclass{article}
\usepackage[czech]{babel}
\usepackage{luavlna}
\usepackage{lua-widow-control}
\usepackage{fontspec}
\usepackage{hyperref}
\usepackage[]{listings}

% šipka pro zalomení dlouhých řádků v listings, měla by se skrýt při kopírování z PDF
% \newcommand{\postbreaksymbol}{%
%   \tagmcbegin{tag=Span,actualtext={}}%
%   \texttt{↪ }%
%   \tagmcend
% }

% ale zdá se, že interní příkazy z latexu na tohle zatím nefungujou, takže použijeme balíček accsup
\usepackage{accsupp}

\newcommand{\postbreaksymbol}{%
  \BeginAccSupp{ActualText={}}%
  \texttt{↪ }%
  \EndAccSupp{}%
}


% breaklines a breakatwhitespace jsou důležitý, protože bez nich se nám text nevejde do viewportu a bude se zobrazovat jenom část stránky
\lstset{language=bash, basicstyle=\ttfamily\small, breaklines=true, breakatwhitespace=true, postbreak=\mbox{\postbreaksymbol}}
% \setmainfont{Droid Sans}
\setmainfont{Literata}
\setmonofont{JetBrains Mono}

% tohle je špatná velikost, protože xteink x4 má 220 ppi takže musíme rozměry přepočítat
% \usepackage[paperwidth=480bp,paperheight=800bp,margin=10bp]{geometry}

% tohle je lepší, protože 480bp je 2.18 palce, což je 55.4mm, což je 157.1bp, takže to odpovídá rozměrům displeje
% 480/220 = 2.18, takže 480bp je 2.18 palce, což je 55.4mm, což je 157.1bp
\usepackage[paperwidth=157.1bp, paperheight=261.8bp, margin=5bp]{geometry}

% % Crosspoint má omezenej viewport, na 464×765 pixelů, takže můžeme použít následující deklaraci, ale vypadá to, že
% je stejně lepší použít celou velikost displeje
% \usepackage[
%   paperwidth=151.85bp,
%   paperheight=250.36bp,
%   margin=5bp
% ]{geometry}

\usepackage{linebreaker}

\usepackage[characters=45 ,scale=heptatonic]{responsive}
\usepackage{microtype}
\usepackage{lipsum}
\begin{document}

\linebreakersetup{maxemergencystretch=0.8em}


\section{Jak je třeba vytvořit Epub pro Xteink}

Zatím žádnej firmware pro Xteink X4 nepodporuje PDF, takže je třeba vytvořit soubor, kterej půjde přečíst.
V zásadě máme dvě možnosti -- vytvořit XTC soubor, nebo Epub obsahující stránky PDF vyrenderovaný do PNG. 
XTC je bitmapovej formát, buď černobílej, nebo se 4 odstínama šedi. Existujou různý konvertory, třeba Xtctools pro 
příkazovou řádku, nebo další, který jsou online. Soubor z Xtctools se mi ale ve firmwaru Crosspoint neotevře,
hlásí to chybu, takže je třeba použít online služby. Výsledek je ale stejně vyrenderovanej v nízký kvalitě, takže 
je špatně čitelnej.

Takže další možnost je vygenerovat ten Epub. Pro tvorbu obrázků stránek můžeme použít Pdftocairo:

\begin{lstlisting}
$ pdftocairo -gray -png -scale-to 800 sample.pdf
\end{lstlisting}

Pro parameter \texttt{-scale-to} je třeba zvolit hodnotu, která odpovídá výšce displeje Xteink X4, teda 800 pixelů.
Výsledek pak převádím skriptíkem, co každou stránku obalí do HTML a vytvoří z nich Epub. Pořád to vypadá trošku hůř,
než jak se zobrazí normální textovej Epub, což je trošku opruz. Ale aspoň je vidět, že to jde.

V Crosspointu je nastavenej omezen viewport kvůli statusbaru, takže ve skutečnosti je viewport omezenej na 765 pixelů:

\begin{lstlisting}
$ pdftocairo  -gray -png -scale-to 765 sample.pdf
\end{lstlisting}


\begin{lstlisting}
$ texlua img2epub.lua adresar output.epub
\end{lstlisting}

Ten skriptík vezme všechny PNG soubory v adresáři a udělá z nich Epub. Výsledek pořád neni ideální, protože Crosspoint očividně
ty obrázky dál zmenšuje a nějak převádí, takže výsledek je hůř čitelnej, než byl originál.

Taky jsem zkoušel vytvořit mono soubor z \texttt{pdftocairo} pomocí:

\begin{lstlisting}
$ pdftocairo  -gray -png -scale-to 800 sample.pdf
\end{lstlisting}

Ale výslednej Epub Crosspoint ani neotevřel, skončilo to chybou. 

Další pokus byl vytvořit PNG soubory s 16 barvama, což by mělo bejt množství, co podporuje ten displej. Buď s ditheringem:

\begin{lstlisting}
$ magick sample-*.png -colorspace Gray -dither zhoufang -colors 16 output.png
\end{lstlisting}

Nebo bez ditheringu:

\begin{lstlisting}
$ magick sample-*.png -colorspace Gray -colors 16 output.png
\end{lstlisting}

Bez ditheringu na počítači vypadá výsledek líp a i na čtečce asi vypadá trošku líp, ale stejně je to pořád deformovaný. 

Asi bude třeba vyzkoušet nějaký jiný firmware, ale skoro všechno je založený na Crosspointu, takže to nevypadá moc nadějně.

% \ResponsiveSetup{characters=35,scale=heptatonic}

% \setsizes{35}

% Následující text bude menší, protože zkouším 35 znaků na řádek

\subsection{Další pokus s XTC}

Vyzkoušel jsem ještě jednou vytvořit XTC soubor pomocí \texttt{cbz2xtc} (\protect\url{https://github.com/srokl/cbz2xtc}), 
který vytváří soubor, co přečte i Crosspoint. Teda aspoň ten černobílej, umí i soubor XTCH, se 4 odstínama šedi. 
Ten Crosspoint neotevřel s hláškou \texttt{Memory error}, ale když jsem to zkoušel pozdějec v Crosspoint 1.3, už se otevřel. 
Výsledek je dobrej, písmo není kostičkovaný a stránky se zobrazujou rychle (narozdíl od Papyrusu).

Musel jsem trošku zasáhnout do zdrojáků toho nástroje, protože vždycky rozsekává stránky na několik obrázků. S touhle změnou 
pak můžu spustit: 

\begin{lstlisting}
$ python3 nosplit.py --clean --contrast-boost 2   --no-split --no-dither ~/ebooks/xtctest/sample.pdf
\end{lstlisting}

To je moje modifikovaná verze skriptu \texttt{cbz2xtc.py}. Volbu \texttt{--no-split} jsem přidal já. Contrast boost trošku
zesílí písmo. Dithering je zakázanej, protože pro text vlastně vůbec nefunguje dobře.

Výsledek má občas trošku deformovaný písmenka, ale jinak to vypadá minimálně stejně dobře, jako ten Epub, možná i trošku líp.

\texttt{cbz2xtc} podporuje další volbu, \texttt{--downscale}, která hodně ovlivňuje způsob, jak se ten obrázek rozpixeluje. Podporuje následující volby:

\begin{description}
  \item[bicubic] -- to je defaultní metoda, průměrný výsledky
  \item[bilinear] -- lepší volba
  \item[box] -- nepoužitelný, rozpixelovaný
  \item[lanczos] -- asi nejlepší
  \item[nearest] -- nefunguje dobře, rozpixelovaný
\end{description}

Tahle volba zatím dává nejlepší výsledky i pro menší písmo (zkoušel jsem to s volbou \texttt{characters=45} balíčku Responsive):
\begin{lstlisting}
$ python nosplit.py --downscale bicubic --2bit --contrast-boost 2 --no-dither  --no-split sample.pdf 
\end{lstlisting}

% \ResponsiveSetup{characters=40,scale=heptatonic}

% \setsizes{40}

% Následující text bude ještě menší, protože zkouším 40 znaků na řádek

\subsection{Pokus s Papyrus}

Papyrus je alternativní firmware, co vypadá, že má širší podporu pro Epub a XTC. Má dost odlišný ovládání než Crosspoint nebo 
původní firmware, takže je třeba se s ním trochu seznámit. Po nahrání 2bitových XTC souboru, co jsem vytvořil pomocí \texttt{cbz2xtc}, 
se mi ho podařilo otevřít, což je pozitivní. Rendering taky vypadá dobře, i když při přechodu na každou stranu se nejdřív 
zobrazí v nižším rozlišení a teprv po zlomku vteřiny se zobrazí vyhlazenej text. Obecně je ten firmware pomalej a otevření 
souboru trvá několik sekund, ale aspoň se to zobrazuje dobře.

Problém je u Epubu. Tam trvá citelně dlouho přechod na každou novou stránku a navíc zobrazuje prázdnou stránku po každym obrázku. 
Obrázky samotný nerenderuje o nic líp, než Crosspoint. Taky v README píšou, že podporujou BMP soubory v Epubu, což je teda 
nevalidní, ale i tak jsem to zkusil. A nefunguje to. Naděje byla, že šestnácti barvou bitmapu to vykreslí s většíma detailama, 
než Crosspoint.

Závěr je, že XTCH se zobrauje dobře, ale je to obecně pomalý, tak budu ještě zkoušet dál.

% input{cviceni}
% \input{epigram}
\end{document}
